\clearpage
\begin{examplebox}
\textbf{例 1：先找直径端点，少做一次配方。}
已知 $A(-1,0),B(1,0)$，$PA/PB=2$，求轨迹。
在直线 $AB$ 上，内分点 $D$ 满足 $AD:DB=2:1$，故 $D=(1/3,0)$；
外分点 $E$ 在 $B$ 的外侧，$AE=2BE$，由 $AE=AB+BE$ 得 $BE=2$，故 $E=(3,0)$。
由前面的直径图，圆心为两端点中点，半径为端点距离的一半：
\[
 C=(5/3,0),\quad R=4/3,\qquad
 \boxed{(x-5/3)^2+y^2=16/9}.
\]
\textbf{回验：} $D$ 处距离为 $4/3,2/3$；$E$ 处为 $4,2$，比值均为 $2$。
证明一已保证整圆周都满足，不能仅凭两个点的回验判定全部轨迹。

\tcbbreak
\textbf{例 2：一般坐标，无须先平移或旋转。}
已知 $A(0,0),B(4,2)$，$PA/PB=2$，求轨迹。
\begin{center}\begin{tikzpicture}[font=\small]
\draw[black!30,rounded corners=2pt] (0,0) rectangle (14.4,8.2);
\node[anchor=west,font=\small\bfseries] at (.5,7.72) {定点不必水平：圆心仍在直线 $AB$ 上，无须先旋转坐标};
\begin{scope}[shift={(3,1.7)},scale=0.8]
\coordinate (A) at (0.0000000000,0.0000000000);
\coordinate (B) at (4.0000000000,2.0000000000);
\coordinate (C) at (5.3333333333,2.6666666667);
\coordinate (P) at (5.3333333333,5.6480906367);
\draw[blue,very thick] (C) circle (2.9814239700);\draw[black!45,dashed] (A)--(C);\draw[blue,thick] (A)--(P);\draw[teal!80!black,thick] (B)--(P);\draw[orange!85!black,dotted,thick] (C)--(P);\fill[black] (A) circle (1.5pt) node[below left=6pt] {$A$};
\fill[black] (B) circle (1.5pt) node[below=7pt] {$B$};
\fill[black] (C) circle (1.5pt) node[below right=14pt] {$C$};
\fill[black] (P) circle (1.5pt) node[above=6pt] {$P$};

\end{scope}
\node[anchor=west] at (.8,6.7) {$A(0,0),\ B(4,2),\ PA/PB=2$};\node[anchor=west] at (.8,.6) {$C(16/3,8/3),\quad R=4\sqrt5/3$};
\end{tikzpicture}\end{center}

因为 $AB=\sqrt{4^2+2^2}=2\sqrt5$，一般公式直接给出
\[
 C=\left(\frac{0-4\times4}{1-4},\frac{0-4\times2}{1-4}\right)
   =(16/3,8/3),\qquad
 R=\frac{2\times2\sqrt5}{3}=\frac{4\sqrt5}{3}.
\]
所以 $\boxed{(x-16/3)^2+(y-8/3)^2=80/9}$。
\textbf{独立回验：} 距离平方条件为
\[
 x^2+y^2=4\bigl((x-4)^2+(y-2)^2\bigr)
 \Longleftrightarrow 3x^2+3y^2-32x-16y+80=0,
\]
配方与答案一致。图中 $A,B,C$ 共线，圆心并不是原点或 $AB$ 的中点。

\tcbbreak
\textbf{例 3：附加直线，求的是交点。}
已知 $A(-3,0),B(3,0)$，$PA/PB=2$，且 $P$ 在直线 $y=x$ 上，求 $P$。
\begin{center}\begin{tikzpicture}[font=\small]
\draw[black!30,rounded corners=2pt] (0,0) rectangle (14.4,8.2);
\node[anchor=west,font=\small\bfseries] at (.5,7.72) {轨迹圆与附加直线联立：整圆缩为两个满足条件的点};
\begin{scope}[shift={(4,3.3)},scale=0.62]
\coordinate (A) at (-3.0000000000,0.0000000000);
\coordinate (B) at (3.0000000000,0.0000000000);
\coordinate (C) at (5.0000000000,0.0000000000);
\coordinate (U) at (1.1771243445,1.1771243445);
\coordinate (V) at (3.8228756555,3.8228756555);
\draw[blue,very thick] (C) circle (4);\draw[black!45] (-3.6,0)--(9.5,0);\draw[purple,thick,dashed] (-2,-2)--(5.6,5.6);\fill[black] (A) circle (1.5pt) node[below=5pt] {$A$};
\fill[black] (B) circle (1.5pt) node[below=5pt] {$B$};
\fill[black] (C) circle (1.5pt) node[below=5pt] {$C$};
\fill[black] (U) circle (1.5pt) node[above left=7pt] {$P_-$};
\fill[black] (V) circle (1.5pt) node[above left=7pt] {$P_+$};

\end{scope}
\node[anchor=west] at (.8,6.8) {圆：$(x-5)^2+y^2=16$};\node[anchor=west,text=purple] at (.8,6) {直线：$y=x$};\node[anchor=west] at (.8,.6) {只取 $P_-,P_+$ 两点};
\end{tikzpicture}\end{center}

先由 $a=3,\lambda=2$ 得 $C=(5,0),R=4$。再代入 $y=x$：
\[
 (x-5)^2+x^2=16
 \Longleftrightarrow 2x^2-10x+9=0
 \Longleftrightarrow x=\frac{5\pm\sqrt7}{2}.
\]
故 $\boxed{P_\pm=((5\pm\sqrt7)/2,(5\pm\sqrt7)/2)}$。
\textbf{回验：} 两点均满足 $y=x$；且
\[
 PA^2-4PB^2=-3\bigl((x-5)^2+y^2-16\bigr)=0.
\]
两点都不是 $B$，距离取正，原比值为 $2$。
判别式 $100-72=28>0$，与图中两个交点一致，不能只取正号解。


\tcbbreak
\textbf{例 4（补充）：距离比小于 $1$，怎样不画错外分点？}
已知 $A(-3,0),B(3,0)$，$PA/PB=1/2$，求轨迹。
\textbf{先定位：} 因为 $PA<PB$，点应更靠近 $A$，轨迹圆在 $A$ 一侧。
在直线 $AB$ 上，内分点 $D$ 满足 $AD:DB=1:2$，由 $AB=6$ 得
$AD=2$，所以 $D=(-1,0)$。
外分点 $E$ 应在 $A$ 左侧：$BE=AB+AE=6+AE$，
而 $AE/BE=1/2$，于是 $2AE=6+AE$，$AE=6$，故 $E=(-9,0)$。
\[
 C=\left(\frac{-1-9}{2},0\right)=(-5,0),\qquad
 R=\frac{|{-9}-({-1})|}{2}=4.
\]
所以
\[
 \boxed{(x+5)^2+y^2=16}.
\]
\textbf{选一个非轴上点核验：} $P(-5,4)$ 在圆上，且
$PA=\sqrt{(-5+3)^2+4^2}=2\sqrt5$，
$PB=\sqrt{(-5-3)^2+4^2}=4\sqrt5$，
故 $PA/PB=1/2$。整圆轨迹的充分性仍由证明一保证。
这个例子与例 1 的 $\lambda=2$ 相互对照，可防止误认为
\textbf{外分点永远在 $B$ 的右侧}。
\end{examplebox}
